\section{An isolated multiple and the integral basis}\label{sec:basis}

Fix a genus $g$ and a multiplicity $m$. We will compare two
representatives associated with these integers. On the equivariant
side the curve and its balanced bundle are fixed throughout the
construction. On the geometric side the entire holomorphic neighborhood
may depend on $m$.

\begin{lemma}\label{lem:germ}
There is a fixed complex threefold germ about a genus-$g$ curve $C$
with acyclic normal bundle in which every $j[C]$, $1\le j\le m$, is
isolated. Its unit $J_m$ is an integral geometric local curve element of
type $(g,m,0)$, and $J_1=x_{g,1}$.
\end{lemma}
\begin{proof}
For $g\ge2$ take the divisorial deformation of a bundle neighborhood
constructed in the proof of \cite[Lemma~6.7]{Pardon}. Its embedded curve
is unobstructed and has normal bundle of index zero, so it has no
infinitesimal deformations in the fixed threefold. The $m$-multiple lies
in the interior semi-regular locus where the semi-chart of multiples of
deformations of $C$ is open \cite[Lemma~6.3 and Definition~6.6]{Pardon}.
At Chern number zero its fixed-target dimension is zero. After restriction
to a tube retracting onto $C$, every component has Chern number zero.
The point $m[C]$ is therefore open in the entire cycle space of the
germ. The inverse image of that point under addition of $(m-j)[C]$ is
$\{j[C]\}$, proving the lower-degree assertions. The complex point
orientation gives the geometric generator. Deformation to the normal
bundle preserves the regular embedded curve in degree one. Its local
chart is a section over the deformation parameter by the implicit
function theorem. On the normal bundle the unit in degree one is
the class with empty evaluation tuple, namely $x_{g,1}$.

In genus zero take $\cO(-1)^{\oplus2}$. In genus one take
$L\oplus L^{-1}$ with $L$ of infinite order in $\mathrm{Pic}^0(C)$.
A nonzero fiber coordinate on a finite cover is impossible by negative
degree in genus zero. In genus one it would trivialize a pulled-back
line, whose norm would make $L$ torsion. In both cases every compact
curve is supported on the zero section, and $J_j=a_{g,j}=x_{g,j}$ by
the empty-tuple construction.
\end{proof}

\begin{proposition}\label{prop:tube}
For a germ in Lemma~\ref{lem:germ} of genus $g\ge2$, there exist integers
$n_{e,h}$ of finite support, $1\le e\le m$, for which
\begin{equation}\label{eq:tube}
 1+\sum_{j=1}^m J_jq^j =
 \prod_{e,h} A_h(q^e)^{n_{e,h}}\pmod {q^{m+1}},
 \qquad n_{1,h}=\delta_{gh}.
\end{equation}
\end{proposition}
\begin{proof}
We place a small tube of the germ in a closed symplectic six-torus.
In a product of two symplectic elliptic curves take one horizontal torus
and $g-1$ distinct vertical translates. Smoothing the $g-1$ transverse
positive intersections gives a symplectic genus-$g$ surface. Its class
pairs once with a vertical torus, hence stays primitive in the product
with a third elliptic curve. Denote the resulting class by $A$ and take
an integral symplectic form. An orientation-preserving identification
of the surfaces may be used to identify their complex tangent structures
along the core by choosing an appropriate compatible complex structure
on its symplectic tangent plane. Their complex normal bundles are
isomorphic, since both have rank two and degree $2g-2$. Choose a tubular
identification with this normal derivative, taking the complex normal
complement to the symplectic orthogonal. The transported integrable
structure is compatible along the core and tame on a small tube.
Extending a compatible structure along the surface and then using a
cutoff in a chart of tame structures realizes the given germ inside
the closed target.

Set $b=\omega(A)>0$ and bound energy by $mb$. There is a tube $V$
with compact closure in this realization such that the only bounded
nonempty cycles in $\overline V$ are $j[C]$ for $j\le m$. Otherwise
a sequence of tubes shrinking to $C$ contains counterexamples. Retraction
to $C$ and positivity of area give class $jA$ with $1\le j\le m$.
Take a subsequence of fixed class. Bounded-area cycle compactness
\cite[Theorem~2.11]{DIW} forces a limit $j[C]$. Convergence of currents
and supports in the integrable tube contradicts isolation in the
holomorphic cycle space. Inside such an isolating tube take $V$ with
compact closure and with $C$ in its interior. Every bounded cycle
contained in $\overline V$ lies in the larger isolating tube, so none
can meet $\partial V$. This property persists under small
perturbations by the same compactness argument. Containment in $V$
therefore defines a compact open and closed locus over a short closed
path to a generic tame structure. Its unit transports $\Phi(J_j)$
in each class $jA$.

At the generic endpoint, the bounded connected supports are disjoint
embedded curves. Partition the supports in $V$ into finitely many
clusters, including the exterior as a separate part
\cite[Propositions~3.18 and~3.21]{DIW}. In each part the union of
the underlying curves is compact by closed incidence. Distinct parts have
disjoint unions, and hence disjoint open target neighborhoods. A
disconnected cycle then decomposes uniquely into cycles in the respective
neighborhoods, including empty cycles. The unit for the union is the
product of the corresponding series, with the convolution performed in
integral homology. This is the excision argument of \cite[\S4]{WueProduct}.
When a path is considered in the ambient threefold, we keep only the
finitely many labels occurring below the energy cutoff. Indeed the
bounded cycle space over the closed path is compact and the integral
homology label is locally constant. Products and inverses of these
series still have finite support at the cutoff, since their positive
degrees have a positive minimum.

For one of these clusters, of class $eA$, apply the cluster isotopy theorem
\cite[Theorem~3.25]{DIW} in the geometric form of
\cite[\S4, proof of Proposition~4.1]{WueProduct}. On a regular arc a
subdivision changes the chosen factor only by factors in proper
multiples. At a fold, compact continuation of the outer insertion gives
the product of the two inner factors equal to one modulo such corrections.
Consequently a crossing inverts the distinguished factor, modulo a
finite product of corrections. At the elementary endpoint its factor is
$\Phi(A_h(q^e))$. Iterate on the correction clusters; their core classes
are $keA$, $k>1$, and $\lfloor m/(ke)\rfloor<\lfloor m/e\rfloor$.
This gives a finite signed product in $S$ at the required cutoff. The
use of disjoint target neighborhoods in excision is what allows these
geometric identities to include the cycles with several components.

Addition and division for the isolated units satisfy
\eqref{eq:localops} up to $m$: the decompositions of $j[C]$ are precisely
$i[C]+(j-i)[C]$. The product on the right of \eqref{eq:tube} obeys
the same laws because substitution, multiplication and inversion
preserve them. Lemma~\ref{lem:detect} lifts the identity from $S$ to $H$.
In degree one it gives $x_{g,1}=\sum_h n_{1,h}x_{h,1}$, so independence
makes the asserted leading exponent positive and equal to one.
\end{proof}

\begin{proof}[Proof of Theorem~\ref{thm:basis}]
Take the coefficient of $q^m$ in \eqref{eq:tube}. The unique unshifted
factor is $A_g(q)$. Any other factor starts in degree $e\ge2$ and only
uses coefficients of multiplicities at most $\lfloor m/e\rfloor<m$.
The binomial coefficients in all integer powers, including negative
powers, are integers. Hence $J_m=a_{g,m}+P_{g,m}$ with $P_{g,m}$
an integral polynomial in lower multiplicities. Choose one germ for
each pair $(g,m)$ as in Lemma~\ref{lem:germ}. Pardon's integral
generation theorem permits \emph{any one such choice for each type}
\cite[Theorem~6.13]{Pardon}. Nonnegative Chern numbers exclude other
types from this summand. Induction in $m$, simultaneously for all $g$,
identifies the subalgebra of $H$ generated by these integral images with
$\ZZ[a_{g,m}]$.

To see that it is a polynomial ring, apply the least-degree relation
argument of Lemma~\ref{lem:detect} to formal balanced generators.
Their logarithmic coordinates are primitive and have the same division
law, and their degree-one images are again the independent $x_{g,1}$.
There is no rational polynomial relation. The triangular equations with
leading coefficient one imply that the chosen $J_m$ for all types
also have algebraically independent images over $\QQ$. Since those
classes generate the original integral summand, its rationalization map
has zero kernel. This proves the stated integral identification.
\end{proof}
